\subsection{2.5 Der rohe OPP als hochdimensionale Struktur}

Die erste Phase der numerischen Untersuchungen verfolgte ein bewusst einfaches Ziel:

\begin{center}
\emph{Welche Eigenschaften besitzt der OPP ohne jede Optimierung und ohne jede Projektion?}
\end{center}

Diese Frage ist von grundlegender Bedeutung.

Denn wenn der OPP bereits intrinsisch eine vierdimensionale Raumzeit enthalten würde, wäre die gesamte Projektionstheorie überflüssig.

Die Ergebnisse der Simulationen zeigten jedoch ein völlig anderes Bild.

%%%%%%%%%%%%%%%%%%%%%%%%%%%%%%%%%%%%%%%%%%%%%%%%%%%%%%%%%%%%%%

\section*{Der Aufbau des rohen OPP}

Der OPP wurde als Menge relational gekoppelter Attraktoren modelliert.

Es existieren keine kartesischen Koordinaten einer Raumzeit.

Vielmehr besitzt jedes Attraktorpaar lediglich eine Relation

\[
R_{ij}.
\]

Diese Relation wurde zunächst durch einen exponentiellen Kernel beschrieben

\[
R_{ij}
=
\exp(-d_{ij}),
\]

wobei

\[
d_{ij}
\]

eine projektive Distanz zwischen zwei Attraktoren bezeichnet.

Aus dieser Matrix wurde anschließend die projektive Gram-Matrix konstruiert.

Deren Eigenwertspektrum bildete die Grundlage sämtlicher Analysen.

%%%%%%%%%%%%%%%%%%%%%%%%%%%%%%%%%%%%%%%%%%%%%%%%%%%%%%%%%%%%%%

\section*{Messgröße: Die effektive Dimension}

Als zentrales Diagnosewerkzeug wurde die effektive Dimension verwendet

\[
d_{\mathrm{eff}}
=
\frac{\left(\sum_i |\lambda_i|\right)^2}
{\sum_i \lambda_i^2},
\]

wobei

\[
\lambda_i
\]

die Eigenwerte der Gram-Matrix darstellen.

Dieses Maß beschreibt nicht die topologische Dimension,

sondern die effektive Zahl dominanter Freiheitsgrade.

%%%%%%%%%%%%%%%%%%%%%%%%%%%%%%%%%%%%%%%%%%%%%%%%%%%%%%%%%%%%%%

\section*{Erste Beobachtung}

Bereits die ersten Simulationen ergaben überraschend konsistente Werte.

Typischerweise erhielt man

\[
30
<
d_{\mathrm{eff}}
<
45.
\]

Die genaue Zahl hing zwar von

\begin{itemize}

\item der Punktzahl,

\item dem Kernel,

\item der Stichprobe

\end{itemize}

ab,

bewegte sich jedoch stets in diesem Größenbereich.

Dies war das erste robuste numerische Resultat der PST.

%%%%%%%%%%%%%%%%%%%%%%%%%%%%%%%%%%%%%%%%%%%%%%%%%%%%%%%%%%%%%%

\section*{Skalierungsversuche}

Um auszuschließen,

dass die hohe Dimension lediglich ein Endlichkeitseffekt ist,

wurden verschiedene Punktzahlen untersucht.

Typische Ergebnisse waren beispielsweise

\[
\begin{array}{c|c}
N & d_{\mathrm{eff}}\\
\hline
10 & \approx 9.6\\
20 & \approx17.8\\
40 & \approx31.6
\end{array}
\]

Die effektive Dimension wächst also mit zunehmender Zahl der Attraktoren.

Dieses Verhalten spricht deutlich gegen einen intrinsischen vierdimensionalen Raum.

Stattdessen deutet es auf einen hochdimensionalen Informationsraum hin.

%%%%%%%%%%%%%%%%%%%%%%%%%%%%%%%%%%%%%%%%%%%%%%%%%%%%%%%%%%%%%%

\section*{Interpretation}

Dieses Ergebnis war zunächst enttäuschend.

Die ursprüngliche Hoffnung lautete,

dass der OPP spontan eine Struktur mit

\[
d_{\mathrm{eff}}
\approx
4
\]

entwickeln würde.

Stattdessen zeigte sich:

\[
\boxed{
\text{Der rohe OPP besitzt keine erkennbare Raumzeitstruktur.}
}
\]

Gerade diese negative Aussage besitzt jedoch hohen wissenschaftlichen Wert.

Denn sie schließt eine ganze Klasse möglicher Modelle aus.

%%%%%%%%%%%%%%%%%%%%%%%%%%%%%%%%%%%%%%%%%%%%%%%%%%%%%%%%%%%%%%

\section*{Robustheit}

Besonders bemerkenswert war,

dass sich die hohe effektive Dimension gegenüber vielen Änderungen als ausgesprochen stabil erwies.

Unter anderem wurden variiert

\begin{itemize}

\item Punktzahl,

\item Zufallsinitialisierung,

\item Stichproben,

\item Kernelparameter.

\end{itemize}

Trotzdem blieb die effektive Dimension stets hoch.

Dies deutet darauf hin,

dass der OPP keine zufällige Punktwolke,

sondern eine eigene hochdimensionale Struktur besitzt.

%%%%%%%%%%%%%%%%%%%%%%%%%%%%%%%%%%%%%%%%%%%%%%%%%%%%%%%%%%%%%%

\section*{Ein wichtiger Fehlschluss wurde vermieden}

An dieser Stelle wäre es leicht gewesen,

durch geeignete Parameterwahl künstlich

\[
d_{\mathrm{eff}}
=
4
\]

zu erzwingen.

Bewusst wurde genau dies nicht getan.

Die numerischen Ergebnisse wurden akzeptiert,

obwohl sie den ursprünglichen Erwartungen widersprachen.

Gerade dadurch blieb das Forschungsprogramm wissenschaftlich belastbar.

%%%%%%%%%%%%%%%%%%%%%%%%%%%%%%%%%%%%%%%%%%%%%%%%%%%%%%%%%%%%%%

\section*{Die erste große Hypothese}

Aus den Ergebnissen entstand erstmals eine neue Vermutung.

Vielleicht gilt gar nicht

\[
\boxed{
\text{Raumzeit ist eine Eigenschaft des OPP.}
}
\]

sondern vielmehr

\[
\boxed{
\text{Raumzeit ist eine Eigenschaft der Beobachtung des OPP.}
}
\]

Dies war der eigentliche Beginn der Projektionstheorie.

%%%%%%%%%%%%%%%%%%%%%%%%%%%%%%%%%%%%%%%%%%%%%%%%%%%%%%%%%%%%%%

\section*{Konsequenzen für die PST}

Die Ergebnisse führten zu mehreren unmittelbaren Schlussfolgerungen.

\begin{enumerate}

\item
Der OPP ist wesentlich hochdimensional.

\item
Die hohe Dimension ist stabil.

\item
Vierdimensionale Raumzeit entsteht nicht automatisch.

\item
Es muss einen zusätzlichen Mechanismus geben.

\end{enumerate}

Damit rückte zwangsläufig die Frage nach geeigneten Projektionsoperatoren in den Mittelpunkt.

%%%%%%%%%%%%%%%%%%%%%%%%%%%%%%%%%%%%%%%%%%%%%%%%%%%%%%%%%%%%%%

\section*{Der Blickwechsel}

Bis zu diesem Zeitpunkt lautete die Forschungsfrage

\[
\text{„Wie erzeugt der OPP Raumzeit?“}
\]

Nach den ersten Simulationen wurde daraus

\[
\boxed{
\text{„Welche Beobachtung macht aus dem OPP eine Raumzeit?“}
}
\]

Dieser Perspektivwechsel gehört zu den wichtigsten konzeptionellen Fortschritten der gesamten PST.

%%%%%%%%%%%%%%%%%%%%%%%%%%%%%%%%%%%%%%%%%%%%%%%%%%%%%%%%%%%%%%

\section*{Zusammenfassung}

Die erste numerische Phase der PST führte zu einem klaren Ergebnis.

\begin{itemize}

\item
Der rohe OPP besitzt eine stabile hochdimensionale Informationsstruktur.

\item
Eine intrinsische vierdimensionale Raumzeit konnte nicht gefunden werden.

\item
Die hohe Dimension ist keine numerische Instabilität.

\item
Die Suche muss sich daher von der Dynamik auf geeignete Projektionen verlagern.

\end{itemize}

Gerade das Scheitern der ursprünglichen Erwartung erwies sich damit als entscheidender wissenschaftlicher Fortschritt.

Die folgenden Kapitel beschäftigen sich deshalb nicht mehr mit dem OPP allein,

sondern mit den mathematischen Projektionen, durch welche eine beobachtbare Physik entstehen könnte.

